% 底本 EGAIV.1.pdf，上传PDF第113--115页，原分册印刷第114--116页。
\subsection{相对形式光滑代数}
\source{第四卷第一部分，印刷第114--115页；上传PDF第113--114页}
\egas{19.9.1}{定义}
设 $\Lambda$ 为拓扑环，$A$ 为拓扑 $\Lambda$-代数，$B$ 为拓扑 $A$-代数。若对每个离散拓扑 $A$-代数 $C$、其中的每个幂零理想 $\mathfrak j$ 及每个连续 $A$-同态 $u_0:B\to C/\mathfrak j$，只要存在连续 $\Lambda$-同态 $u:B\to C$ 使
\[
 u_0:B\xrightarrow{u}C\longrightarrow C/\mathfrak j,
\]
就也存在连续 $A$-同态 $v:B\to C$ 使
\[
 u_0:B\xrightarrow{v}C\longrightarrow C/\mathfrak j,
\]
则称 $B$ 为\emph{相对于 $\Lambda$ 形式光滑的 $A$-代数}。

由定义，若 $B$ 是形式光滑 $A$-代数，则对 $A$ 上每一种拓扑 $\Lambda$-代数结构，即每个连续环同态 $\Lambda\to A$，$B$ 都相对于 $\Lambda$ 形式光滑。

\egas{19.9.2}{命题}
设 $\Lambda$ 为拓扑环，$A$ 为拓扑 $\Lambda$-代数。

$(i)$ $A$ 作为 $\Lambda$-代数，相对于 $\Lambda$ 形式光滑。

$(ii)$ 若 $B$ 为相对于 $\Lambda$ 形式光滑的 $A$-代数，$C$ 为相对于 $\Lambda$ 形式光滑的 $B$-代数，则 $C$ 为相对于 $\Lambda$ 形式光滑的 $A$-代数。

$(iii)$ 若 $B$ 为相对于 $\Lambda$ 形式光滑的 $A$-代数，$A'$ 为拓扑 $A$-代数，则拓扑 $A'$-代数 $B\otimes_AA'$ 相对于 $\Lambda$ 形式光滑。

$(iv)$ 设 $B$ 为拓扑 $A$-代数，$S\subset A$、$T\subset B$ 为乘法集，且 $S$ 的典范像包含于 $T$。若 $B$ 为相对于 $\Lambda$ 形式光滑的 $A$-代数，则 $T^{-1}B$ 为相对于 $\Lambda$ 形式光滑的 $S^{-1}A$-代数。

$(v)$ 设 $B_1,\ldots,B_n$ 为拓扑 $A$-代数。则 $\prod_{i=1}^nB_i$ 相对于 $\Lambda$ 形式光滑，当且仅当每个 $B_i$ 都如此。

$(i)$ 直接由定义得到；其余证明紧密仿照（19.3.5），原著留给读者。
\cmt{H8-N04}{这里的“相对”要求已知一个 $\Lambda$-线性提升，再使其成为 $A$-线性提升；不能省去这个前提。特别是 $(i)$ 不声称任意 $\Lambda$-代数在通常意义下形式光滑。$(v)$ 应采用本译稿（19.3.5，v）已经修补的正交幂等元提升论证，不能沿用初版对乘积源同态的错误分解。}

\egas{19.9.3}{推论}
设 $\Lambda$ 为拓扑环，$A',B$ 为两个拓扑 $\Lambda$-代数。则拓扑 $\Lambda$-代数 $B'=B\otimes_\Lambda A'$，作为 $B$-代数相对于 $\Lambda$ 形式光滑。

由（19.9.2，i、iii）得到。

\egas{19.9.4}{命题}
设 $\Lambda$ 为拓扑环，$A$ 为拓扑 $\Lambda$-代数，$B$ 为拓扑 $A$-代数。以下条件等价：

$a)$ $B$ 为相对于 $\Lambda$ 形式光滑的 $A$-代数；

$b)$ $\widehat B$ 为相对于 $\Lambda$ 形式光滑的 $A$-代数；

$c)$ $\widehat B$ 为相对于 $\Lambda$ 形式光滑的 $\widehat A$-代数；

$d)$ $\widehat B$ 为相对于 $\widehat\Lambda$ 形式光滑的 $\widehat A$-代数。

证明仿照（19.3.6），原著仍留给读者。

\egap{19.9.5}
同样，把（19.3.8）中的“形式光滑”替换为“相对于 $\Lambda$ 形式光滑”，结论及证明仍有效。对（19.3.10）作同样替换时，结论应改为：每个连续 $A$-同态 $u_0:B\to C/\mathfrak j$，若可通过连续 $\Lambda$-同态 $u:B\to C$ 提升，则也可通过连续 $A$-同态 $v:B\to C$ 提升。证明基本不变。

\egap{19.9.6}
把“形式光滑”替换为“相对于 $\Lambda$ 形式光滑”，（19.4.1）和（19.4.2）的判据仍成立，证明几乎不变。

\egas{19.9.7}{命题}
设 $\Lambda$ 为拓扑环，$A$ 为拓扑 $\Lambda$-代数，$B$ 为拓扑 $A$-代数。假设对每个离散 $A$-代数 $C$、每个满足 $\mathfrak j^2=0$ 的理想 $\mathfrak j\subset C$，每个连续 $A$-同态 $u_0:B\to C/\mathfrak j$，若可提升为连续 $\Lambda$-同态 $B\to C$，则也可提升为连续 $A$-同态 $B\to C$。那么 $B$ 为相对于 $\Lambda$ 形式光滑的 $A$-代数。

（19.4.3）的证明可直接用于此处。

\egas{19.9.8}{命题}
设 $\Lambda$ 为拓扑环，$A$ 为拓扑 $\Lambda$-代数，$B$ 为拓扑 $A$-代数。则 $B$ 相对于 $\Lambda$ 形式光滑，当且仅当对每个被 $B$ 的某个开理想零化的离散拓扑 $B$-模 $L$，有
\[
 \Exalcotop_{A/\Lambda}(B,L)=0.
\]
这里沿用（18.4.2）及（18.5.2）的相对扩张群记号。

沿用（19.4.4）证明的记号，只须注意此时可以假设 $B_\lambda$ 的扩张 $E_\lambda$ 为 $\Lambda$-平凡扩张，证明的其余部分不变。

当 $\Lambda,A,B$ 都离散时，此判据化为
\begin{equation}
 \Exalcom_{A/\Lambda}(B,L)=0
 \quad\text{对每个 $B$-模 $L$ 成立}.
 \tag{19.9.8.1}
\end{equation}
换言之，$B$ 的每个以 $B$-模为核的交换 $A$-扩张，只要为 $\Lambda$-平凡扩张，就也是 $A$-平凡扩张。

\subsection{形式非分歧代数与形式平展代数}
\source{第四卷第一部分，印刷第115--116页；上传PDF第114--115页}
\egap{19.10.1}
我们还将使用两个与形式光滑代数相近的概念。

\egas{19.10.2}{定义}
设 $A$ 为拓扑环，$B$ 为拓扑 $A$-代数。若对离散 $A$-代数之间每个核幂零的 $A$-同态 $p:E\to C$ 及每个连续 $A$-同态 $u:B\to C$，至多存在一个连续 $A$-同态 $v:B\to E$，使 $u=p\circ v$，则称 $B$ 为\emph{形式非分歧 $A$-代数}。

若 $B$ 同时形式光滑且形式非分歧，称其为\emph{形式平展（formellement étale）$A$-代数}。等价地，对每个核幂零的满 $A$-同态 $p:E\to C$，其中 $E,C$ 离散，以及每个连续 $A$-同态 $u:B\to C$，存在唯一连续 $A$-同态 $v:B\to E$，使 $u=p\circ v$。

\egas{19.10.3}{例}
这里只讨论离散环。

$(i)$ 若结构同态 $\rho:A\to B$ 满射，则 $B$ 为形式非分歧 $A$-代数：在（19.10.2）的记号下，复合 $A\xrightarrow{\rho}B\xrightarrow{v}E$ 必为给定结构同态，所以 $v$ 若存在则唯一。

令 $I=\Ker\rho$。此时，$B$ 形式光滑，当且仅当 $I=I^2$；这一条件也等价于 $B$ 形式平展。
\cmt{H8-E01}{初版此处说“反之，$B$ 只有在 $\rho$ 双射时才形式光滑”，一般不成立。取任一非零域 $k$，投影 $\rho:k\times k\to k$ 满射而不单射，其核 $0\times k$ 满足 $I=I^2$，故形式平展。正文给出准确判据及译者补证，不将其冒称为原作者已刊勘误。}
事实上，必要性可取平方零扩张 $A/I^2\to A/I=B$，并提升 $B$ 的恒等同态。提升若为 $A$-线性，必使 $I$ 在 $A/I^2$ 中的像为零，所以 $I\subseteq I^2$，从而 $I=I^2$。反之，给定离散 $A$-代数 $E$、幂零理想 $J\subset E$，以及 $A$-同态 $B\to E/J$，设结构同态为 $\alpha:A\to E$，则 $\alpha(I)\subseteq J$。若 $J^N=0$，由 $I=I^N$ 得
\[
 \alpha(I)=\alpha(I^N)\subseteq J^N=0.
\]
因此 $\alpha$ 唯一分解为 $A\to B\to E$，给出所需提升。

$(ii)$ 若 $S$ 为环 $A$ 的乘法集，则 $B=S^{-1}A$ 为形式平展 $A$-代数。设 $j:A\to B$ 为典范同态，$\rho:A\to E$ 为环同态，$J\subset E$ 幂零，$p:E\to C=E/J$ 为商同态。若 $u:B\to C$ 为 $A$-同态，则对每个 $s\in S$，$u(s/1)=p(\rho(s))$ 可逆。由于 $J$ 幂零，$\rho(s)$ 本身也可逆（$0_{\mathrm I}$，7.1.12），所以 $\rho$ 唯一分解为 $A\xrightarrow{j}B\xrightarrow{v}E$。由 $u\circ j=(p\circ v)\circ j$ 及局部化的泛性质，$u=p\circ v$。

$(iii)$ 域 $k$ 的有限可分扩张 $k'$ 为形式光滑 $k$-代数（19.6.1）。后面将证明：对 $k$ 的有限生成扩张 $K$，形式平展 与形式非分歧等价，而且也等价于 $K/k$ 为有限可分扩张（21.7.4）。

$(iv)$ 多项式代数 $A[X_\alpha]_{\alpha\in I}$ 对 $A$ 形式光滑（19.3.3）。当 $A\ne0$ 且 $I\ne\varnothing$ 时，它不形式非分歧，因而不形式平展。
\cmt{H8-T03}{初版只排除 $I=\varnothing$，未排除零基环 $A=0$；零环上的多项式代数仍是零环，结构同态为恒等，形式平展。正文写明非零条件。若 $A\ne0$，选一个变量分别映为 $0$ 和 $\varepsilon\in A[\varepsilon]/(\varepsilon^2)$，其余变量映为零，即得同一模 $\varepsilon$ 同态的两个不同提升。}

\egas{19.10.4}{注}
与（19.4.3）相同的论证表明，验证形式非分歧时，可以在（19.10.2）中只考虑 $\Ker p$ 平方为零的情形。还可以只考虑 $u(B)\subseteq p(E)$ 的情形，否则不存在提升；再以 $u(B)$ 代替 $C$，以 $p^{-1}(u(B))$ 代替 $E$，即可进一步假设 $u$ 和 $p$ 都满射。
